% ============================================================
% 67_code_comments.tex — OPTIONALES Code-Kommentar-Modul
% ============================================================
% Opt-in: in preamble.tex NACH 65_code_style einkommentieren.
% Smart Inline / Overflow Platzierung annotierter Code-Zeilen
% ausserhalb von lstlisting (in contentbox / defbox).
% Verantwortung dieser Datei ist nur die Platzierung; Längen,
% Farben und Schrift werden aus den Style-Modulen bezogen:
%
%   \ZSFcodeCommentMaxWidth        -> 30_layout_spacing
%   \ZSFcodeCommentThreshold       -> 30_layout_spacing
%   \ZSFcodeCommentGap             -> 30_layout_spacing
%   \ZSFcodeCommentMaxWidthDefault -> 30_layout_spacing (immutabel)
%   \ZSFcodeCommentThresholdDefault-> 30_layout_spacing (immutabel)
%   \ZSFfontCodeComment            -> 50_typography_semantics
%   ce_green                       -> 40_colors_structure
%
% Public API:
%   \CodeLine{code}[comment]      Auto: Kommentar rechts, sonst oben
%   \CodeLine{code}               Nur Code
%   \InlineComment{text}          Kommentar rechts (kein Auto-Layout)
%   \OverlineComment{text}        Kommentar in eigener Zeile davor
%   \SetCodeCommentThreshold{len} Schwelle anpassen (persistent bis Reset)
%   \ResetCodeCommentThreshold    Schwelle auf Default zurücksetzen
%
% Hinweise:
% - Funktioniert NICHT in lstlisting; in lstlisting native '#'
%   Kommentare benutzen.
% - Standard-Aufrufer-Pattern in Kapiteln endet jede gerenderte
%   Zeile mit '\\' (ttfamily-Block), daher fuegt diese Datei kein
%   abschliessendes '\\' an, ausser bei \OverlineComment, wo der
%   Zeilenwechsel Teil der Semantik ist.
%
% Herkunft: konsolidiert aus ZSF_Informatik2_v2/styles/67_code_comments.tex.
% ============================================================

\makeatletter

% ------------------------------------------------------------
% Sicherheits-Guard: setzt voraus dass 30_layout_spacing.tex
% bereits geladen wurde (Reihenfolge in preamble.tex).
% ------------------------------------------------------------
\@ifundefined{ZSFcodeCommentThreshold}{%
  \PackageError{ZSF/code-comments}%
    {Required register \string\ZSFcodeCommentThreshold\space missing}%
    {Load styles/30_layout_spacing.tex before styles/67_code_comments.tex.}%
}{}
\@ifundefined{ZSFfontCodeComment}{%
  \PackageError{ZSF/code-comments}%
    {Required macro \string\ZSFfontCodeComment\space missing}%
    {Load styles/50_typography_semantics.tex before styles/67_code_comments.tex.}%
}{}

% ------------------------------------------------------------
% Render-Helfer
% ------------------------------------------------------------
% Code wird verbatim-aehnlich gerendert: \detokenize macht alle
% Tokens cat-12, sodass '#', '_', '%', '&' im Code-Argument
% erscheinen wie geschrieben.
\newcommand{\ZSF@CC@code}[1]{\texttt{\detokenize{#1}}}
% Kommentare bekommen automatisch ein '# '-Praefix (Python-Stil) und
% erben die Code-Schriftgröße (kein \small), damit Code und Kommentar
% optisch auf einer Zeile sitzen.
\newcommand{\ZSF@CC@comment}[1]{{\ZSFfontCodeComment\#\,#1}}

% Mess-Speicher für Auto-Layout
\newsavebox{\ZSF@CC@codebox}
\newsavebox{\ZSF@CC@cmtbox}
\newlength{\ZSF@CC@codewd}
\newlength{\ZSF@CC@cmtwd}
\newlength{\ZSF@CC@effectiveThreshold}

% ------------------------------------------------------------
% \CodeLine{code} oder \CodeLine{code}[comment]
% ------------------------------------------------------------
% \NewDocumentCommand mit "m o"-Spec verwendet xparse-konforme
% Argument-Erkennung; vermeidet \@ifnextchar[ und damit chktex-
% Bracket-Counting-Probleme.
\NewDocumentCommand{\CodeLine}{m o}{%
  \IfNoValueTF{#2}%
    {\ZSF@CC@code{#1}}%
    {\ZSF@CC@autoPlace{#1}{#2}}%
}

\newcommand{\ZSF@CC@autoPlace}[2]{%
  \sbox{\ZSF@CC@codebox}{\ZSF@CC@code{#1}}%
  \sbox{\ZSF@CC@cmtbox}{\ZSF@CC@comment{#2}}%
  \setlength{\ZSF@CC@codewd}{\wd\ZSF@CC@codebox}%
  \setlength{\ZSF@CC@cmtwd}{\wd\ZSF@CC@cmtbox}%
  % Effektive Schwelle: min(ZSFcodeCommentThreshold, \columnwidth) -- wird
  % zur Render-Zeit ausgewertet, sodass der korrekte Spaltenbreite-Wert
  % aus der aktiven multicols-Umgebung greift.
  \setlength{\ZSF@CC@effectiveThreshold}{\ZSFcodeCommentThreshold}%
  \ifdim\ZSF@CC@effectiveThreshold>\columnwidth
    \setlength{\ZSF@CC@effectiveThreshold}{\columnwidth}%
  \fi
  \ifdim\dimexpr\ZSF@CC@codewd+\ZSF@CC@cmtwd+\ZSFcodeCommentGap\relax
        > \ZSF@CC@effectiveThreshold\relax
    \ZSF@CC@renderAbove
  \else
    \ZSF@CC@renderRight
  \fi
}

% Code links, Kommentar rechts mit Mindest-Gap.
\newcommand{\ZSF@CC@renderRight}{%
  \leavevmode
  \usebox{\ZSF@CC@codebox}%
  \hspace*{\ZSFcodeCommentGap}%
  \usebox{\ZSF@CC@cmtbox}%
}

% Kommentar in eigener Zeile davor, dann Code.
\newcommand{\ZSF@CC@renderAbove}{%
  \leavevmode
  \usebox{\ZSF@CC@cmtbox}%
  \newline
  \usebox{\ZSF@CC@codebox}%
}

% ------------------------------------------------------------
% \InlineComment{text} -- immer rechts, kein Auto-Layout.
% Erwartet, dass davor bereits Code im selben Zeilenkontext
% steht; fuegt nur den Kommentar mit Mindest-Gap an.
% ------------------------------------------------------------
\newcommand{\InlineComment}[1]{%
  \hspace*{\ZSFcodeCommentGap}\ZSF@CC@comment{#1}%
}

% ------------------------------------------------------------
% \OverlineComment{text} -- Kommentar in eigener Zeile, gefolgt
% von einem Zeilenumbruch, sodass die naechste \CodeLine-Zeile
% auf der Folgezeile beginnt.
% ------------------------------------------------------------
\newcommand{\OverlineComment}[1]{%
  \leavevmode\ZSF@CC@comment{#1}\newline
}

% ------------------------------------------------------------
% Threshold-API
% ------------------------------------------------------------
% Persistent: bleibt bis zum naechsten Reset oder Gruppen-Ende.
\newcommand{\SetCodeCommentThreshold}[1]{%
  \setlength{\ZSFcodeCommentThreshold}{#1}%
}

\newcommand{\ResetCodeCommentThreshold}{%
  \setlength{\ZSFcodeCommentMaxWidth}{\ZSFcodeCommentMaxWidthDefault}%
  \setlength{\ZSFcodeCommentThreshold}{\ZSFcodeCommentThresholdDefault}%
}

\makeatother
